% ch-output.tex -- what a run writes and how to read it
\chapter{Output files}
\label{ch:output}

Every run writes into one directory, named after \key{JobName}. This chapter
says what each file contains, in what units, and which of them you should look
at when something goes wrong. The keys that control them are in
\cref{sec:ref-bcout:output}.

\section{The result directory}
\label{sec:output:dir}

\begin{yzntable}{Contents of a result directory\label{tab:resultdir}}{@{}L{0.30\textwidth} L{0.62\textwidth}@{}}
\textsf{\bfseries File} & \textsf{\bfseries Contents}\\ \midrule
\file{<input file name>}          & a verbatim copy of the input file that produced the run \\
\file{JobInfo.txt}                & version, build, machine, start and end time, and the linear-solver backend of each equation \\
\file{Warnings.txt}               & warnings raised during the run; zero-length if there were none \\
\file{FatalErrors.txt}            & why the run aborted; zero-length if it did not, see \cref{sec:output:warnings} \\
\file{field00000.vts}, \dots\     & VTK structured grid, one per field output \\
\file{field00000.dat}, \dots\     & the same field as tab-separated text; also the ASCII restart format \\
\file{<JobName>.pvd}              & ParaView time series referencing the \file{.vts} files \\
\file{History.dat}                & integral quantities and extrema against time \\
\file{CalcInfo.dat}               & solver diagnostics against time \\
\file{Probe00.dat}, \dots\        & time series at the points listed in \key{Output.ProbeList} \\
\end{yzntable}

The input file is copied in deliberately: a result directory carries the
complete description of the run that made it, so it stays interpretable after
you have edited the case twelve more times.

One thing you may expect here is not here: the scales, the dimensionless
groups, the progress log and the termination message are written to
\textbf{standard output}, not to a file in this directory
(\cref{sec:output:running}). Redirect if you want to keep them.

\key{JobName} is a path relative to the working directory, and the
directories on the way to it are made as needed: \val{"runs/sweep/case1"}
writes into \file{runs/sweep/case1/}, and the \file{.pvd} in it is named after
the last part, \file{case1.pvd}. Spaces and non-ASCII characters (Japanese,
say) are allowed in it, as in the name of the input file and in
\key{Restart.FileName}.

Before anything is computed the run makes sure the directory can be used, and
stops with exit code 1 and a message naming the path and \key{JobName} when it
cannot: when it already exists and \key{Output.Overwrite} is \val{false}, when
something other than a directory has its name, when it cannot be made (a
parent that is a file, no permission) and when no file can be written in it
(a read-only directory, no permission). Nothing is written then, and the
message is on standard error only, since \file{FatalErrors.txt} would be in
the directory that failed.

\begin{note}[Two runs can share a directory]
  The directory is named after \key{JobName}, not after the input file. Two
  input files with the same \key{JobName} write to the same place, and with
  \key{Output.Overwrite} set to \val{true} the second replaces the files of
  the first that it writes again: \file{History.dat}, \file{CalcInfo.dat}, the
  probes, the \file{.pvd}, \file{JobInfo.txt} and the field files of the
  numbers it reaches. Field files it does not reach --- \file{field00012.*}
  of a first run longer than the second --- are left as they were, and the
  new \file{.pvd} does not list them. This is convenient for a parameter sweep
  run in separate working directories and a trap for one run in the same
  directory; remove the directory to start clean.
\end{note}

\section{\file{JobInfo.txt}}
\label{sec:output:jobinfo}

Six sections: program info (version, build time, compiler, the versions of the
input reader, grid and utility layers), solver info, machine info (OS, host,
CPU, processor count, the OpenMP thread count \code{OMP\_MAX\_THREADS} the run
used, the count the runtime offered \code{OMP\_THREADS\_REQUESTED} and the
cells per thread of the thread limit \code{OMP\_THREAD\_CAP}
(\cref{sec:parallel:cap}), the algorithms in force \code{YUZEN\_OPT}
(\cref{tab:envvars}), the kind of OpenMP runtime \code{OMP\_RUNTIME} and the
grain of the linear solves \code{YUZEN\_GRAIN} (\cref{sec:parallel:runtime}),
memory, process id), starting
info (date, time, working directory, result directory), performance (the
wall-clock time of the time-step loop alone, \code{Time loop [s]}, the steps
taken, the Newton and linear iterations (the Newton count is 0 in \val{"CZ"},
which has no Newton loop, \cref{sec:output:calcinfo}), in the well-mixed model the steps
failed because the Newton correction stagnated far from convergence
(\cref{sec:convergence:stagnation}), in \val{"CZ"} the coupling passes
taken again and the steps failed because they did not settle
(\cref{sec:verification:czpasses}), the largest Courant number of each
limited convection scheme, in a well-mixed run with a concentration the
solute moved from above $c = 1$ and the solute cut there, in either model
the solute moved into cells below $c = 0$ and the solute made up there
(\cref{sec:ref-numerics:convection}), and the sum of the steps solved
against the times the loop started and ended at) and ending info (elapsed
time, date, time, and the termination type). The termination type is
\code{normal} for a run that ended on one of the stops of
\cref{tab:termination}, \code{interrupted (SIGINT)} or
\code{interrupted (SIGTERM)} for one stopped by a signal
(\cref{sec:output:failures}) and \code{abort} for one stopped by a fatal
error. A run killed outright (\code{SIGKILL}, a power cut) has no ending info
at all: the file ends after the starting info.

\code{\#\#\#\#\#\# Solver Info \#\#\#\#\#\#} is the section that depends on
your input. It names the model, and the linear-solver backend and the
convection scheme that each equation actually used:

\begin{plain}
###### Solver Info ######
Model :                      Well-Mixed-1D
Linear solver <height> :     eigen
Linear solver <conc> :       eigen
Convection scheme <height> : Koren
Convection scheme <conc> :   UTOPIA
\end{plain}

The \code{<conc>} lines appear only when the model solves a concentration
field, and the \val{"CZ"} model, which takes the flux of the film height
centered, has no \code{Convection scheme <height>} line. \textbf{These are resolved names, not the strings you wrote}: with
\key{Backend} at its default \val{"auto"} the line reports whichever
implementation was chosen, and after the silent \val{"amgcl"} fall-back
(\cref{sec:ref-numerics:linearsolver}) it reports \val{"eigen"}, the one that
ran. The backend lines (and, in the well-mixed model, the scheme lines) are
echoed to standard output under the heading
\code{\#\#\#\#\#\# Solver \#\#\#\#\#}. \Cref{ch:linearsolver} compares the
backends.

Everything else in the file describes the program, the machine and the wall
clock. The rest of what the solver computed from your input --- the scales, the
dimensionless groups, the grid --- goes to standard output
(\cref{sec:output:running}).

\section{The run log on standard output}
\label{sec:output:running}

Everything the solver has to say while it runs goes to \textbf{standard
output}: it is a terminal log, not a file, and no result directory contains it.
Keep it by redirecting:

\begin{shell}
yuzen case.toml > run.log
\end{shell}

\file{tools/run\_case.sh} already does this, capturing standard output and
standard error together into \file{stdout.txt} inside the run directory
(\cref{sec:env:tools}), which is the easiest way to end up with a log you can
read later without thinking about it.

This is where the solver reports what it built from your input, and then how
the run progressed. It opens with the linear-solver backend and the convection
scheme of each equation and the grid, then the scales
used to non-dimensionalize the problem ($L_0$, $H_0$, $T_0$, $U_0$ under the
heading \code{Scaling factors}), then the dimensionless groups $\As$, $\Ma$,
$\Bo$, $\Pgx$, $\Pgy$, $\Spn$, $\Dpn$, $\Pen$, $\Df$, $\Ev$ and $\Mr$ under
\code{Nondimensional numbers}, then the non-dimensional time and output
intervals, and the model name.

Read the scales here when a result is wrong by orders of magnitude rather than
by a few percent: compare $L_0$, $T_0$ and $U_0$ against what you expected
before looking anywhere else, because a mis-set
\key{Fluid.\bk SurfaceTension.\bk BaseValue} or
\key{Initial.\bk SurfaceShape.\bk TotalHeight} shows up here immediately and
nowhere else for a long time. \Cref{sec:governing:scales} defines each scale
and \cref{sec:governing:groups} each group.

Below that block comes the progress log, written as the run goes:

\begin{plain}
.....................Field file 4, Time = 1.285130e-05
.............................dt+ (1.013153e-07)
..........
############## Time reached to the EndTime. #############
Field file 5, Time = 1.606806e-05
Largest Courant number of Numerics.Convection.Height = "Koren": 0.000469

############## Program normally terminated. #############
\end{plain}

One dot is one time step. \code{dt+} and \code{dt-} mark the adaptive step
growing and shrinking, with the new value in parentheses. \code{Field file
$n$} marks a field written, with its non-dimensional time.

The line that says why the loop ended is the one that matters; after it come
the field of the last step, the largest Courant number of each limited
convection scheme (\cref{sec:ref-numerics:convection}) and the closing line.
There are three normal terminations:

\begin{yzntable}{Termination messages\label{tab:termination}}{@{}L{0.33\textwidth} L{0.59\textwidth}@{}}
\textsf{\bfseries Message} & \textsf{\bfseries Meaning}\\ \midrule
\code{Time reached to the EndTime.} & The run reached \key{Numerics.EndTime}. \\
\code{Dry out.}                     & The concentration test selected by \key{Physics.\bk Evaporation.\bk EndConcType} passed \key{EndConcValue}. \\
\code{Steady state.}                & The change rate stayed below \key{Numerics.\bk SteadyStateCriterion} for \key{SteadyStateSteps} consecutive steps. \\
\end{yzntable}

A run stopped by Ctrl-C or \code{SIGTERM} ends instead with
\code{Program stopped by SIGINT; its output is complete.} (or
\code{SIGTERM}), and a run that diverged with \code{Program terminated after
divergence.}; \cref{sec:output:failures} lists the exit codes.

A long alternation of \code{dt-} and \code{dt+} is worth investigating: it
usually means the time step has locked onto the level at which the Newton
convergence test is just satisfiable, which is a convergence problem and not a
time-stepping one (\cref{sec:convergence:absolute}). A run that reports
\code{Steady state.} after a handful of steps on a fine mesh is likewise
suspicious; \cref{sec:convergence:termination} explains why, and why the
criterion is a rate rather than a per-step change.

\section{\file{Warnings.txt} and \file{FatalErrors.txt}}
\label{sec:output:warnings}

Both files are opened as soon as the result directory is created, and
\textbf{neither carries a banner}. A zero-length \file{Warnings.txt} means no warning
was raised, and a zero-length \file{FatalErrors.txt} means the run did not
abort. Testing the two files for emptiness is a valid check on a finished run,
and it is the cheapest one there is.

The warnings you are most likely to meet are the Newton stagnation notice
(\cref{sec:convergence:stagnation}), the linear-solver iteration-limit notice, and
the notice that \val{"amgcl"} was requested in a build without it. A limited
convection scheme run past a Courant number of one adds two lines, grep for
\code{Courant number} (\cref{sec:ref-numerics:convection}); so does a
well-mixed film that holds more solute than liquid and has some cut at
$c = 1$, grep for \code{cut at c = 1}, and a film with less than no solute
that has some made up at $c = 0$, grep for \code{made up at c = 0}. None of
them stops the run. The stagnation notice is written once, for the first
step cut for it; \file{JobInfo.txt} counts them all.

\begin{note}[Three messages, verbatim, and what to grep for]
  Three messages are worth having in full, because a finished run is usually
  diagnosed by grepping \file{Warnings.txt} or \file{FatalErrors.txt} for one
  of them. A short distinctive substring is the right thing to search for, and
  one is given with each:

  \begin{itemize}
    \item \textbf{Newton stagnation}
      (\file{numerics/\bk NewtonKantorovich.cpp}): the line gives the iteration
      the correction stagnated at, the correction reached and the first one,
      and says the step is cut, as in
      \code{Newton-Kantorovich: the Newton correction stagnated far from
      convergence (iteration 7, correction 6.586796e-01 against 4.827437e-01
      in the first); the step is cut and solved again (JobInfo.txt counts the
      steps cut for this; later ones are not reported here).} At a fixed step
      the same cause is the fatal error that stops the run. Grep for
      \code{correction stagnated}.
    \item \textbf{The \val{"amgcl"} fall-back}
      (\file{numerics/\bk LinearSolverBackend.cpp}):
      \code{LinearSolver: Backend="amgcl" was requested, but this build has no
      AMGCL. Falling back to "eigen".} Grep for \code{has no AMGCL}.
    \item \textbf{The precursor-ratio abort}
      (\file{model/\bk NonDimensionalNumbers.cpp}), which is fatal and goes to
      \file{FatalErrors.txt}: \code{Physics.DisjoiningPressure requires
      Initial.SurfaceShape.PrecursorRatio > 0.} Grep for
      \code{PrecursorRatio}, which appears verbatim.
  \end{itemize}
\end{note}

A fatal error does. It is written through a tee (\code{TeeStream} in
\file{yuzen/src/util/YuzenUtil.h}) that duplicates it into
\file{FatalErrors.txt} \emph{and} standard error, so the message reaches the
terminal and the result directory at once. The file holds the reason and then
the abort marker:

\begin{plain}
TimeStepAdjustment(). Under flow: the requested step fell below Numerics.TimeStepControl.MinDt.

####### Aborting ############
\end{plain}

The conditions that end a run this way are: disjoining pressure enabled with
\key{Initial.\bk SurfaceShape.\bk PrecursorRatio} left at \val{0}
(\file{model/\bk NonDimensionalNumbers.cpp}); the adaptive step falling below
\key{Numerics.\bk TimeStepControl.\bk MinDt}
(\file{numerics/\bk ImplicitSolver.cpp}); a step that does not converge
with the time-step control off, so that it cannot be retried at a shorter
step --- the well-mixed Newton iteration, or a linear solve of the film or
of the concentration, the message naming which
(\file{numerics/\bk NewtonKantorovich.cpp},
\file{model/\bk ConcZ/\bk CZ2D/\bk CZ\_HeightEquation\_2D.cpp}); a
\val{"Pinning"} edge meeting a \val{"FreshInlet"} or an \val{"Outlet"} edge
at a 2-D corner, which is not implemented
(\file{model/\bk BoundaryCondition.cpp}); a zero denominator in the Laplace
pressure kernel (\file{model/\bk well\_mixed/\bk LaplacePressureKernel.h}); and
a model that does not exist, from an unsupported combination of
\key{Physics.\bk ConcField.\bk Model} and \key{Grid.Dimension}
(\file{Yuzen.cpp}); an output file that cannot be written
(\cref{sec:output:failures}); and running out of memory, which names the
grid as the likely cause (\file{mainProg.cpp}).

\begin{warn}[An input-file error leaves no result directory at all]
  \code{ReadInputFile()} runs before \code{CreateDirAndFiles()}, so an input
  that fails the consistency check --- or a \key{SolverName} the parser
  rejects --- aborts with the message on \textbf{standard error only}. The
  same holds for a command line the solver does not accept, an input file
  that is missing, a directory, empty or not readable, a result directory
  that cannot be used (\cref{sec:output:dir}) and a \key{Restart.FileName}
  that is not there. No result directory is created, and there is therefore
  no \file{FatalErrors.txt} to look in. If you go hunting for one after a
  rejected input and find nothing, this is why: read the terminal, or the
  \file{stdout.txt} of \file{tools/run\_case.sh} (\cref{sec:env:tools}).
\end{warn}

\section{A full disk, an interrupted run, and the exit code}
\label{sec:output:failures}

\paragraph{Every file is whole or absent.} A field file (\file{.vts},
\file{field*.dat}), its \file{field*.state} and the \file{.pvd} are written
under their name with \file{.tmp} added and renamed when the last byte is on
the disk. Whatever stops the run while one is being written --- a full disk,
a quota, a file size limit, a kill --- leaves at worst a \file{*.tmp} beside
the files, which nothing reads; a file under its real name is never cut
short. The \file{.pvd} is rewritten whole, closing tags included, after every
field file, so ParaView opens it while the run is going on and after it was
killed. \file{History.dat}, \file{CalcInfo.dat} and the probes are appended
one row at a time; a run killed in the middle of a row leaves that row without
its newline, and a restart into the same directory drops it.

\paragraph{A file that cannot be written stops the run.} When a write fails ---
no space left on the device, a file grown past a limit
(\code{ulimit -f}), the directory made read-only under the run --- the run
stops with exit code 1, after naming the file and the reason on standard
error and in \file{FatalErrors.txt}:

\begin{plain}
Could not write run/field00007.vts: No space left on device.
The run is stopped: its output could not be written into run (full disk, ...
\end{plain}

It is not continued without its output, and it does not write the
\file{.state} of a field file that failed. The field files already in the
directory are complete, and the last one that has its \file{.state} beside it
can be restarted from (\cref{sec:ref-numerics:restart}) once there is room.

\paragraph{Ctrl-C and \code{SIGTERM} end the run cleanly.} The first
\code{SIGINT} (Ctrl-C) or \code{SIGTERM} --- what a batch scheduler sends
at the end of a job's time --- lets the step in progress finish, writes the
outputs of the state reached as at the end of a run (a field file and its
\file{.state}, the rows of \file{History.dat}, \file{CalcInfo.dat} and the
probes), records \code{interrupted (SIGINT)} in \file{JobInfo.txt} and ends
with exit code $128$ plus the signal number: 130 for \code{SIGINT}, 143 for
\code{SIGTERM}, which is what the shell reports for a process the signal
killed. The run continued from that field file and its \file{.state} is,
bit for bit, the run that was not interrupted (\code{error\_paths}). A second
Ctrl-C ends the process at once; so does \code{SIGKILL}, and the files are
then still whole --- restart from the last field file that has its
\file{.state} beside it (the newest one lacks it when the run was killed
between writing the two).

\begin{yzntable}{Exit codes\label{tab:exitcodes}}{@{}L{0.12\textwidth} L{0.80\textwidth}@{}}
\textsf{\bfseries Code} & \textsf{\bfseries Meaning}\\ \midrule
0   & The run ended on one of the stops of \cref{tab:termination}; \code{--check-input}, \code{--help} and \code{--version} succeeded. \\
1   & Anything that went wrong: the command line, the input file, the result directory, the restart file, a write, a fatal error during the run, a diverged run, running out of memory. \\
130, 143 & Stopped by \code{SIGINT} or \code{SIGTERM}; the output of the state reached is complete. \\
\end{yzntable}

\section{\file{field*.dat}}
\label{sec:output:field}

Tab-separated text with a one-line header naming each column and its unit
(\code{xx[mm]}, \code{HS[um]}, \code{Ux[m/s]}, \code{CC[-]}, and so on, or
\code{[-]} throughout when \key{Output.Dimensioned} is \val{false}). The
\val{"CZ"} header is additionally prefixed with \code{\#}. The columns depend
on the model:

\begin{yzntable}{Columns of \file{field*.dat}\label{tab:field-columns}}{@{}L{0.16\textwidth} L{0.76\textwidth}@{}}
\textsf{\bfseries Model} & \textsf{\bfseries Header}\\ \midrule
1-D well-mixed & \code{xx}\quad\code{HS}\quad\code{Ux}\quad\code{CC}\quad\code{HH}\quad\code{SS}\quad(\code{PHI}) \\
2-D well-mixed & \code{xx}\quad\code{yy}\quad\code{HS}\quad\code{Ux}\quad\code{Uy}\quad\code{CC}\quad\code{HH}\quad\code{SS}\quad(\code{PHI}) \\
\val{"CZ"}, 1-D & \code{xx}\quad\code{zz}\quad\code{CC}\quad\code{Ux}\quad\code{Uz}\quad\code{HS}\quad\code{HH}\quad\code{SS} \\
\val{"CZ"}, 2-D & \code{xx}\quad\code{yy}\quad\code{zz}\quad\code{CC}\quad\code{Ux}\quad\code{Uy}\quad\code{Uz}\quad\code{HS}\quad\code{HH}\quad\code{SS} \\
\end{yzntable}

Every field file carries a header naming each column and its unit, and every
one of them follows \key{Output.Dimensioned}.

The rows are points, not cells. Along each in-plane axis there is one point on
each wall and one at every cell center, so a 2-D file has $(N_x+2)(N_y+2)$
points per layer; the value on a wall is the mean of the cell next to it and
the ghost cell beyond, and a velocity component there is the one on that face.
A \val{"CZ"} file does the same through the thickness: one layer on the
substrate, one at each of the $N_z$ cell centers, one on the free surface,
$(N_x+2)(N_y+2)(N_z+2)$ rows in 2-D and $(N_x+2)(N_z+2)$ in 1-D. Its top layer
is therefore the concentration at the free surface. The data run $x$ fastest,
then $y$, then $z$.

\code{HS} is the free-surface height $h+s$ and \code{HH} is the film thickness
$h$; over a flat substrate they coincide. \code{Ux} is the \emph{mixed}
velocity, the sum of the contributions of every active term, which is also what
advects the concentration --- \cref{sec:verification:velocity} is the test that
pins that composition rule down.

With \key{Output.Dimensioned} set to \val{true} the units are
\textbf{not all the same}: the in-plane coordinates \code{xx} and \code{yy} are
in \si{\milli\metre}, while \code{HS}, \code{HH} and \code{SS} are in
\si{\micro\metre} and \code{Ux}, \code{Uy} and \code{Uz} are in
\si{\metre\per\second} --- \code{Uz}, the vertical velocity $w$ of a \val{"CZ"}
file from continuity, as well: it is made dimensionless with $H_0/T_0$ rather than $L_0/T_0$,
and is converted back with that scale.
\code{CC} is a mass fraction and has no unit. The mixed in-plane and
out-of-plane scaling is easy to miss when you plot column 2 against column 1:
the aspect ratio of such a plot is exaggerated by a factor of $10^{3}$.

In a \val{"CZ"} file the \code{zz} column is the \textbf{height} of the point,
\begin{equation*}
  z = s + h\,\zeta ,
\end{equation*}
not the mapped coordinate $\zeta \in [0,1]$ the model works in
(\cref{sec:conc:cz}). It therefore runs from \code{SS} to \code{HS} at each
$(x,y)$ and takes a different set of values in each column of the film, which
is the same thing the matching \file{field*.vts} carries in its point
coordinates. Plot \code{CC} against \code{zz} and the vertical axis is
\si{\micro\metre} above the datum, like \code{HS} and \code{SS} beside it.

\subsection{The driving potential}
\label{sec:output:potential}

With \key{Output.\bk FieldData.\bk WritePotential} set to \val{true} (well-mixed
model only), a well-mixed field file has one more column, last, \code{PHI},
and its \file{.vts} one more array, \code{Potential}: the driving potential
\begin{equation*}
  \varphi = \hat\sigma(c)\,\nabla^2(H+S) + \Dpn\,\Pihat(H) - \Bo\,(H+S),
  \qquad \Pihat(H) = \Bigl(\frac{H_*}{H}\Bigr)^{n} - \Bigl(\frac{H_*}{H}\Bigr)^{m},
\end{equation*}
whose gradient drives the flux through the mobility, $\Qvec =
\tfrac{H^3}{3\hat\mu}\nabla\varphi$ plus the body forces and the Marangoni
stress (\cref{ch:governing}): a uniform $\varphi$ carries no flux, which is
the equilibrium condition. $\nabla^2$ is the second difference the capillary
term differences once more at each face, so that the flux of the scheme is
the mobility times the difference of \code{PHI} across the face. The in-plane
gravity, the centrifugal force and the Marangoni stress are not gradients of
it and are left out. Each term is zero when its \key{Physics} block is off.

With \key{Output.Dimensioned} \val{true} \code{PHI} is in pascals, on the
pressure scale $\sigma_0 h_0/L_0^2$:
\begin{equation*}
  \code{PHI} = \frac{\sigma_0 h_0}{L_0^2}\,\varphi
  = \sigma\,\nabla^2(h+s) + \Pi(h) - \rho g\,(h+s),
\end{equation*}
the pressure in the film being $p = -\code{PHI}$ (up to a constant); with
\val{false} it is $\varphi$. A wall row takes the value of the cell next to
it. The restart reads past the column, and the \code{Potential} array is
written after the point coordinates, where a restart from the \file{.vts}
does not reach. \code{analytic} checks it (\code{potential\_1d}).

Values are written with 18 significant digits whatever \key{Output.Dimensioned}
says, because this file is also the ASCII restart input. Six digits is not
enough: passing the height through a second difference divides by $\Delta x^2$,
and on a contact-line mesh at $\Delta x = \SI{0.3}{\micro\metre}$ the output
precision alone shifted the pressure by a few percent and the velocity derived
from it by three orders of magnitude.

\section{\file{History.dat}}
\label{sec:output:history}

One row per \key{Output.\bk GeneralHistory.\bk TimeInterval}, with time in the first
column and then:

\begin{yzntable}{Columns of \file{History.dat}\label{tab:history-columns}}{@{}L{0.22\textwidth} L{0.70\textwidth}@{}}
\textsf{\bfseries Column} & \textsf{\bfseries Meaning}\\ \midrule
\code{Volume}         & total liquid volume \\
\code{ResinVol}       & total solute volume \\
\code{AveC}           & mean concentration \\
\code{EvaporatedVol}  & cumulative volume lost to evaporation \\
\code{OutflowedVol}   & cumulative volume that left through the open edges,
                        \val{"FreshInlet"} and \val{"Outlet"} (negative for
                        what came in; in 1-D with the centrifugal term, also
                        what it flings sideways) \\
\code{OutflowedResin} & cumulative solute that left with it, carried and
                        diffused \\
\code{VolErr}         & volume conservation error \\
\code{ResinVolErr}    & solute conservation error \\
\code{maxHS}, \code{minHS} & extrema of the free surface \\
\code{maxH}, \code{minH}   & extrema of the film thickness \\
\code{maxC}, \code{minC}   & extrema of the concentration \\
\code{SuppliedVol}    & cumulative volume the nozzle supplied \\
\code{SuppliedResin}  & cumulative solute that came with it, at the concentration \key{Physics.Source.ConcType} gives it (last column) \\
\end{yzntable}

With \key{Output.Dimensioned} set to \val{true} the time column is in
\si{\second}, but the lengths and volumes here are in \si{\metre} and
\si{\metre\cubed} --- \emph{not} in the \si{\micro\metre} of the field files.
\file{ProbeNN.dat} follows the field files instead. \Cref{sec:ref-bcout:output} says which
writer applies which scaling. A 1-D film is a strip \SI{1}{\metre} wide in $y$:
its volumes are per meter of width, $\sum h\,\Delta x$, and the header gives
their unit as \si{\metre\cubed\per\metre}. The width does not depend on
\key{Grid.Ymin} and \key{Grid.Ymax}, which a 1-D film does not use, nor on the
length scale, and the two concentration models write the same volume for the
same film. With \key{Output.Dimensioned} set to \val{false} the volumes are
dimensionless, $\sum h\,\Delta x$ in 1-D and $\sum h\,\Delta x\,\Delta y$ in 2-D,
on the scales of \cref{sec:governing:scales}.

\code{VolErr} and \code{ResinVolErr} are the residual of the balance between
what is in the domain, what evaporated, what left through the boundaries and
what the nozzle supplied: $\code{VolErr} = (\code{Volume} + \code{OutflowedVol}
+ \code{EvaporatedVol} - \code{SuppliedVol})/\code{Volume}_0 - 1$ and
$\code{ResinVolErr} = (\code{ResinVol} + \code{OutflowedResin} -
\code{SuppliedResin})/\code{ResinVol}_0 - 1$, with the volume the sum over the
cells. The sums use the time step each step was solved with. By default
(\key{Physics.Source.ConcType = "local"}) the nozzle's
liquid arrives at the concentration already where it lands
(\cref{sec:terms:source}), so \code{SuppliedResin} adds up, over each step,
the supply of every cell times its concentration, the mean of the start and
the end of the step; with a concentration uniform in the plane, as in the
shipped nozzle cases, \code{ResinVolErr} is then \code{VolErr}. With
\val{"supplied"} it arrives at \key{Physics.Source.Conc}, and
\code{SuppliedResin} adds up the supply of every cell times \key{Conc},
$\int c_S S$. Either way that is
the solute the concentration equation itself takes in with the supply
(\cref{sec:disc:wm-conservative}), so on a closed film the solute balances
to rounding, with evaporation or without (\code{conc\_drying},
\cref{sec:verification:analytic}), and with a nozzle in either mode
(\code{nozzle\_solute}, \code{nozzle\_supplied}). The
centrifugal term flings a 1-D film sideways, out through the two long sides of the strip (the
$\partial_y$ of the flux that the 1-D equation keeps,
\cref{sec:terms:centrifugal}); that liquid, and the solute it carries, is
counted in \code{OutflowedVol} and \code{OutflowedResin} as the film-height
equation takes it out, so that a spinning 1-D film between walls balances to
rounding (\code{spin\_balance}, \cref{sec:verification:analytic}). The
solute a convection scheme carries past $c = 1$ is moved to the cells around,
not removed, and a cell it carries below $c = 0$ is filled from the cells
around (\cref{sec:disc:convection}), so neither leaves a trace here
(\code{conc\_overshoot}, \code{conc\_undershoot}). What crosses a
\val{"FreshInlet"} or an \val{"Outlet"} is counted, step by step, from the
face fluxes the step itself put through the edge faces: the volume flux of
the film-height equation and the solute flux of the concentration equation
(the convective flux with the scheme's face value and the diffusive flux), at
the Crank--Nicolson weighting of the step, over the committed steps only, so
that an open film balances to rounding as a closed one does
(\code{inlet\_outlet}, \code{solute\_balance}; the running totals are in the
state file, so a restarted run continues the same balance). Counted from the
velocities and the face values of the field after each step instead, the
balance was found off by the error of the scheme, $6\times10^{-7}$ on 100
cells and falling at order 1.8 in $\Delta x$; corrected.

A \val{"CZ"} run writes \code{Volume}, \code{ResinVol}, \code{AveC}, the
extrema (\code{maxHS} to \code{minC}), and then \code{EvaporatedVol},
\code{OutflowedVol}, \code{OutflowedResin}, \code{VolErr} and
\code{ResinVolErr}, formed as above; it has no nozzle and no open edge, so
\code{OutflowedVol} is the sideways outflow of 1-D alone.
\code{VolErr} and \code{ResinVolErr} are the cheapest continuous check that a case is behaving: they should
stay small and they should not drift. Note that they sit near zero, so a
\emph{relative} comparison of them between two runs amplifies rounding --- this
is why they show the largest relative difference in the backend comparison of
\cref{sec:verification:backends} while every physical quantity agrees to
$10^{-10}$.

\section{\file{Energy.dat}}
\label{sec:output:energy}

With \key{Output.Energy.Enable} set to \val{true} (well-mixed model only) the
run writes \file{Energy.dat}, a row with every row of \file{History.dat}: the
time, then
%
\begin{description}
  \item[\code{FreeEnergy}] the free energy of the film,
    \begin{equation*}
      \mathcal F = \int\Bigl[\tfrac12\hat\sigma\,|\nabla(H+S)|^2
        + \tfrac{\Bo}{2}(H+S)^2 + \Dpn\,w(H)\Bigr]\,\mathrm dA,
      \qquad w(H) = -\int_{H_*}^{H}\Pihat(s)\,\mathrm ds,
    \end{equation*}
    the functional whose variation is the driving potential with its sign
    turned, $\delta\mathcal F/\delta H = -\varphi$
    (\cref{sec:output:potential}). Discretely the gradient term is a sum
    over the faces of $\tfrac12\hat\sigma_f(\Delta(H+S)/\Delta x)^2$, so that
    the derivative of $\mathcal F$ with respect to a cell's $H$ is minus
    \code{PHI} of that cell, exactly.
  \item[\code{Dissipation}] $\mathcal D$, the power the potential does
    against the face fluxes of the scheme, $\sum_i\varphi_i R_i\,\Delta A$
    with $R_i$ the rate the face fluxes give cell $i$. For the capillary
    pressure alone it is $\sum_f D_f(\Delta\varphi/\Delta x)^2\Delta A \ge 0$,
    the integral of $\frac{M}{\hat\mu}|\nabla\varphi|^2$ with the face mobility of
    the scheme, and of slip with \key{Physics.Slip}; it is what the
    discretized equations dissipate, so that $\mathrm d\mathcal F/\mathrm dt =
    -\mathcal D$ holds for them exactly in time.
  \item[\code{Dissipated}] $\int_0^t\mathcal D\,\mathrm dt$, by the trapezoidal rule over the
    steps (the Crank--Nicolson weighting).
  \item[\code{Balance}] the discrete energy balance of the step that
    ended at the row, $(\mathcal F^{n+1} - \mathcal F^n)/\Delta t +
    (\mathcal D^n + \mathcal D^{n+1})/2$: zero to $O(\Delta t^2)$ relative
    to $\mathcal D$ with Crank--Nicolson. The first row has 0. Summed over
    the steps it is $\mathcal F(t) - \mathcal F(0) + \code{Dissipated}$,
    which \code{analytic} checks at second order (\code{energy\_*}).
\end{description}
%
With \key{Output.Dimensioned} \val{true} the energies are in joules (per
meter of width in 1-D, J/m), on the scale
$\sigma_0 h_0^2$ ($\sigma_0 h_0^2/L_0$ in 1-D), and the powers in watts (per
meter) on that over $T_0$. The values have 18 significant digits. A run
continued from a state file continues \code{Dissipated} and gives every row
what the uninterrupted run gives it (\code{restart\_roundtrip}); from a field
file alone \code{Dissipated} starts afresh.

\begin{warn}[$\mathcal F$ falls only for the gradient flows]
  $\mathcal F$ is a Lyapunov function --- it never rises --- when the flux
  is the mobility times $\nabla\varphi$ and nothing else: the capillary
  pressure, the vertical gravity and the disjoining pressure, a uniform
  concentration ($\hat\sigma$ and $\hat\mu$ constant), and \val{"Symmetry"}
  or \val{"Periodic"} edges. The in-plane gravity, the centrifugal force,
  the Marangoni stress, evaporation, the nozzle and open edges put energy
  in or take it out; \file{Energy.dat} is written for them too, and
  \code{Dissipation} is then the power of $\varphi$ against the flux, not a
  dissipation. A \val{"Pinning"} edge holds the film with a ghost that
  follows the cell beside it, which $\mathcal F$ does not account for.
  The three terms share one face mobility
  (\cref{sec:disc:positivity-one}), so for them $\mathcal D$ is a sum of
  squares, the disjoining pressure included, and $\mathcal F$ has risen in
  none of the checks. \code{analytic} checks the balance and the fall
  (\code{energy\_*}).
\end{warn}

\section{\file{CalcInfo.dat}}
\label{sec:output:calcinfo}

The solver's own diagnostics, one row per
\key{Output.\bk CalcInfo.\bk TimeInterval}: \code{NK\_iter} and \code{NK\_resid} (the
Newton outer iteration count and final correction), \code{BCG\_iter} and
\code{BCG\_resid} (the linear solver), \code{deltaT}, and \code{ChangeRate}.
A \val{"CZ"} run has no Newton loop of its own and writes \code{Time},
\code{BCG\_iter}, \code{BCG\_resid}, \code{deltaT} and \code{ChangeRate}.
A well-mixed run whose field breaks down --- a NaN, a concentration outside
$[-0.1, 1.1]$, a film below $-10^{-3}$ of its scale, a negative volume ---
stops with \code{terminated after divergence} and writes no final field, as a
\val{"CZ"} run does.

This is the file to open when a run is slow. \code{NK\_iter} sitting at
\key{NewtonMaxIterations} means the Newton loop is not converging --- go to
\cref{ch:convergence}. \code{BCG\_iter} in the thousands where another backend
needs tens means the preconditioner, not the problem
(\cref{sec:linsolve:choosing}). \code{deltaT} collapsing means either a genuine
stiffness or, far more often, the interaction between an absolute convergence
test and the correction floor (\cref{sec:convergence:absolute}).

\code{ChangeRate} is the relative change of the film thickness per unit time,
in \si{\per\second}. It is what the steady-state test is built on, so look at
this column before choosing \key{Numerics.\bk SteadyStateCriterion} rather than
guessing a threshold and finding out later that the run stopped at step twelve.

\section{VTK output and ParaView}
\label{sec:output:vtk}

\file{field*.vts} are VTK structured-grid files and \file{<JobName>.pvd} is a
collection referencing them with their time stamps. Open the \file{.pvd}: you
get the whole series with the correct physical times, rather than a numbered
sequence ParaView has to guess the timing of. The times are written to 12
significant digits, so field files $10^{-4}$\,s apart stay distinct up to
$t = 10^{7}$\,s. The \file{.pvd} lists every field file and is rewritten whole
after each one (\cref{sec:output:failures}), so its cost grows with the
number of outputs: after 30{,}000 field outputs it is 1.6\,MB, and on a
two-core \SI{2.1}{\giga\hertz} Xeon rewriting it took about
\SI{10}{\milli\second} per output. That matters only for a run that writes
tens of thousands of fields.

A \val{"CZ"} \file{.vts} is a three-dimensional structured grid whose points
sit at the heights of the cells, $z = s + h\,\zeta$ for each cell's own
$\zeta$, ghost layers included --- so the two outermost layers at the top and
the bottom lie just outside the film, and the real cells fill it from substrate
to free surface. \code{Velocity} is given at those points, each component the
mean of the two faces of the cell it is stored on. Besides \code{Conc} and
\code{Velocity} it carries
\code{FilmThickness} and \code{Substrate}, each repeated down its column; these
are what a restart from the \file{.vts} reads the film back from.

A well-mixed \file{.vts} carries \code{Conc}, \code{Height} (the free surface
$h+s$, the \code{HS} of the field file), \code{Substrate} and the mixed
\code{Velocity}, and beside it the velocity of each term on its own:
\code{V\_Ca} (capillary pressure), \code{V\_Ma} (Marangoni), \code{V\_Bo}
(vertical gravity), \code{V\_Pg} (in-plane gravity) and \code{V\_Sp}
(centrifugal force). The disjoining pressure's share is in \code{Velocity}
but has no array of its own. These are the way to see which term is moving the
liquid at a given place (\cref{sec:governing:dimensional}). With \key{WritePotential}, \code{Potential} holds the driving potential
(\cref{sec:output:potential}).

The two output formats are independent: \key{Output.\bk FieldData.\bk VTKFile}
and \key{Output.\bk FieldData.\bk NeutralFile}. Turning off the VTK files saves a
significant amount of disk on long 2-D runs; turning off the text files means
you cannot restart from the result, so it is rarely the right economy.

\begin{warn}[The \file{.vts} files honor \key{Output.Dimensioned}, and mix
  their length units]
  Every \file{.vts} writer applies the key the way \file{field*.dat} does:
  in-plane coordinates in \si{\milli\metre}, heights --- and therefore the
  $z$ coordinate of the grid, which is built from them --- in
  \si{\micro\metre}, velocities in \si{\metre\per\second}.
  \code{Conc} is a mass fraction and is never scaled.

  So a \file{.vts} is exaggerated vertically by $10^{3}$ relative to its true
  proportions. That is usually what you want, since a film is orders of
  magnitude thinner than it is wide and true proportions would render as a
  flat sheet; but do not read a slope or an angle off the picture without
  correcting for it, and set the $z$ scale explicitly in ParaView if you want
  the real aspect ratio. With \key{Output.Dimensioned} set to \val{false},
  everything in the file is non-dimensional and the exaggeration is gone.
\end{warn}

\section{Probes}
\label{sec:output:probes}

Each \key{Output.ProbeList.Probe} block records a time series at one point:
free-surface height, concentration, velocity components, film thickness and
substrate height. The coordinates you give are rounded to the nearest grid
point at start-up. One file is written per probe, named \file{ProbeNN.dat} with
\code{NN} a zero-padded two-digit index \textbf{counting from \code{00}} in the
order the blocks appear in the input file --- so two probes give
\file{Probe00.dat} and \file{Probe01.dat}. Probes are written at
\key{Output.ProbeList.TimeInterval},
so they are the way to get a finely resolved time history at a few locations
without writing the whole field at that rate. No case in \file{cases/} uses
them.
